\chapter{From a confident answer to uncertainty about a lending effect}
\label{closure30:graph-confidence}

A credit committee can ask an assistant to propose a causal diagram, attach
probabilities to its arrows and simulate a change in lending terms. Each step
looks familiar. Together they can create an assurance that none of the steps
has earned: a narrow interval for the financial consequence of a policy.
The issue is not whether a language model can help an analyst. It is whether
the meaning of uncertainty survives the passage from an answer about an
arrow to a distribution over diagrams and then to an intervention effect.

Molero P\'erez's master's thesis supplies a concrete, unusually complete
description of this proposed passage. It combines language-assisted discovery,
retrieval, three measures of answer confidence, probabilistic graphs and
causal normalizing flows. It evaluates pairwise direction judgments on the
T\"ubingen cause--effect pairs and a seven-variable medical application using
Framingham observations. The work is valuable as a workflow comparator:
it asks how a proposed graph changes the eventual intervention calculation.
Its probability assignments and evaluation target nevertheless need separate
justification before a bank could use the resulting intervals.
\citep[Sections 4.1--4.4]{closure30molero}

\section{What the complete workflow actually evaluates}

The thesis compares self-reported confidence, answer-token probabilities and
an information-based score called HallBayes. In its final procedure, the last
score supplies a quantity called risk of hallucination, $r_{ij}$. An accepted
arrow is assigned probability $1-r_{ij}$; a rejected arrow is assigned
probability $r_{ij}$. Three runs are averaged. Each possible arrow is then
drawn independently as a Bernoulli variable, producing 100 candidate graphs.
A topological sort is used to obtain an ordering for a causal flow.
\citep[pp.~44--50, equations (14)--(19), Appendix D.3]{closure30molero}

For the medical example, the reference diagram is an expert-constructed
diagram drawn from the clinical literature. The thesis explicitly recognizes
that the true graph of a complex medical population is generally unknown.
It trains 100 models for each of four arrangements: the reference graph,
the language-derived graphs, the reference ordering without adjacency
restrictions, and random orderings without adjacency restrictions. For each
intervention--outcome pair, it calculates effects from generated populations
and takes the 2.5th and 97.5th percentiles across the 100 models.
The principal comparison is overlap between these intervals and intervals
from the fitted reference-graph models.
\citep[pp.~52, 58--62]{closure30molero}

This design measures agreement with a specified fitted reference. It does
not observe the result of assigning patients, or borrowers, to the proposed
intervention. Its percentiles combine the graph construction, training
variation and simulation used in that exercise. They are not established as
95\% confidence intervals under repeated sampling from the population.
The distinction matters even if the neural models are fitted without error.
A model can reproduce an observational distribution under a stipulated
diagram without establishing that the diagram is the right one.

The reported mean interval-overlap scores are 0.723 for estimated graphs,
0.712 for the reference ordering alone, and 0.301 for random orderings.
The first two differ by 0.011. These are reported source results, not a
replication in this proposal. The table does not supply a paired uncertainty
analysis establishing that the small difference is a general advantage of
language-derived graphs. Nor is the PC comparison a comparison of identical
information: the language procedure receives semantic and retrieved
knowledge, while PC works from observations. The thesis specifies
chi-square independence tests while favouring continuous clinical variables;
the public implementation examined here does not reconstruct that full
comparison. Accordingly, the proposal retains the workflow and its measured
target, without importing a claim of general causal or commercial superiority.
\citep[pp.~54--55, 77--79, Tables 10--11]{closure30molero}

The thesis itself recognizes consequential limitations. It reports that the
HallBayes calibration curve is not sufficiently good, describes the loss of
valid arrows during refinement, and acknowledges the unaddressed possibility
of hidden common causes. Those admissions support a measured reading of its
contribution. The analysis below makes the corresponding probability and
decision requirements explicit rather than attributing an unconditional
guarantee to every part of the source.
\citep[pp.~68--69 and 79]{closure30molero}

\section{A ceiling on success is not a probability of success}

Write $p$ for the probability of a correct answer to a specified binary
question, $q$ for a reference probability, and
\begin{equation}
 d(p,q)=p\log\frac{p}{q}+(1-p)\log\frac{1-p}{1-q},
 \qquad 0<q<1,
 \label{closure30:bernoulli-kl}
\end{equation}
using the usual continuous interpretation at $p=0$ and $p=1$. All logarithms
in this treatment are natural, so information is measured in nats. Suppose,
for the moment, that a valid information bound supplies
$d(p,q)\leq\Delta$. Define
\begin{equation}
 p_{\max}(\Delta,q)=\sup\{u\in[0,1]:d(u,q)\leq\Delta\}.
 \label{closure30:success-cap}
\end{equation}
Then $p\leq p_{\max}$ and hence
\begin{equation}
 1-p\ \geq\ 1-p_{\max}.
 \label{closure30:error-floor}
\end{equation}
The complement is a lower bound on error. It is neither an upper bound on
error nor an estimate of the error probability. On printed page 47, following
equation (16), the thesis calls this complement an upper bound. That direction
is inconsistent with its definition of $p_{\max}$ as the highest feasible
success probability. The inconsistency is visible in the original page;
it is not a text-extraction artefact.
\citep[p.~47, equations (15)--(16)]{closure30molero}

For example, $q=0.1$ and $\Delta=0.5$ give
$p_{\max}=0.49505075$. A success probability of $0.1$ satisfies the inequality
because $d(0.1,0.1)=0$. Its error probability is $0.9$, which is above the
reported complement $0.50494925$. Substituting the cap as the probability of
an arrow would turn a feasible upper limit into a point belief without
additional evidence. The numerical root was checked with both bisection and
an independent bracketing solver.

There is a second issue before even this direction argument applies.
The thesis's equation (14) averages clipped positive log-probability changes
for the selected answer after masking parts of its prompt. The original
information inequality concerns an expectation of divergence between full
distributions, with a specified event and reference. A selected answer's
log-probability change is not that expectation merely because both quantities
are measured in nats. Clipping, selection of the answer and the meaning of
the correctness event all need to be accounted for.
\citep[pp.~46--47]{closure30molero}

\section{What the original information inequality does supply}

The first version of Chlon and colleagues' paper appeared in September 2025
under the title \emph{Predictable Compression Failures: Why Language Models
Actually Hallucinate}. Its June 2026 third version has a narrower title and
an explicit emphasis on evidence-grounded binary adjudication. Both versions
are available before the survey cutoff. Their substantive differences must
be retained: the later version requires a specified verification primitive
or calibrated verifier, and places unrestricted generation without one
outside its stated regime.
\citep{closure30chlonone,closure30chlonthree}

The backbone of Theorem 5 is the following valid inequality. Let $P$ be a
reference distribution, $S_k$ comparison distributions, $\overline S$ their
mixture, and $A$ a fixed event. With $p=P(A)$ and $q=\overline S(A)$,
\begin{equation}
 \underbrace{\E_k\{D_{\rm KL}(P\Vert S_k)\}}_{\Delta}
 \ \geq\ D_{\rm KL}(P\Vert\overline S)
 \ \geq\ d(p,q).
 \label{closure30:information-chain}
\end{equation}
The first inequality follows from convexity in the second argument. The
second follows by reducing the outcome to the indicator of $A$. Equivalently,
partition the likelihood ratio over $A$ and its complement and apply the
log-sum inequality on each part. This supplies a necessary information
requirement for a given event probability. Equality in the second step needs
no remaining distinction within either part. Equality in the full chain also
needs equality in the mixture step; an exponential tilt relative to
$\overline S$ alone does not remove variation among the $S_k$.
\citep[Theorem 5 and Appendix A.4]{closure30chlonone,closure30chlonthree}

A large full-distribution divergence need not increase the event probability.
For a four-point outcome space, take
\begin{equation}
 S=(0.05,0.05,0.001,0.899),\qquad
 P=(0.05,0.05,0.899,0.001),
 \label{closure30:irrelevant-information}
\end{equation}
and let $A$ consist of the first two outcomes. Use the same $S$ in every
comparison, so there is no permutation complication. Both $P(A)$ and $S(A)$
are $0.1$, while
\begin{equation}
 D_{\rm KL}(P\Vert S)=0.898\log(899)=6.10755216,
 \qquad d(0.95,0.1)=1.99420862.
 \label{closure30:budget-counterexample}
\end{equation}
The first quantity exceeds the proposed information requirement for 95\%
success by a factor of about 3.06, yet success remains 10\%. The information
distinguishes outcomes \emph{within failure}; it does not establish success.
This is a project counterexample to treating the necessary full-distribution
budget as a sufficient guarantee. It is not a re-analysis of the paper's
reported experiments.

The third version computes a Bernoulli divergence for a verifier-relative
predicate. This avoids the particular within-failure variation in
\eqref{closure30:irrelevant-information}. Direction still matters:
$d(0.05,0.5)=d(0.95,0.5)$, so a minimum-divergence threshold alone cannot
distinguish 5\% from 95\% success. Restricting $p$ to the increasing branch
$p\geq q$, using an exact binary divergence and knowing what the reference
probability means would remove this particular ambiguity. These are real
conditions, not consequences of a confident answer. A calibrated verifier
would itself require validation for the relevant predicate and population.
\citep[Sections 3.5--3.6, Algorithm 1 and Appendix A.5]
{closure30chlonthree}

For a bank, verifying that a document contains a statement is different from
verifying that the statement identifies a causal arrow. A successful test of
quotation or logical consistency can support the former without supplying
the latter. The practical recommendation is to use such scores as disclosed
diagnostics for evidence review, and to evaluate any decision threshold on a
separate, adjudicated set of the actual questions it will handle. A numerical
threshold should not be passed downstream as a calibrated graph probability
unless that stronger interpretation has been established.

\section{Independent arrows can describe an impossible joint belief}

Even defensible marginal probabilities need not define a distribution over
legal causal diagrams. Consider only the possible arrows
$A\to B$, $B\to C$ and $C\to A$, each assigned probability $4/5$ and sampled
independently. The probability of drawing the complete directed cycle is
\begin{equation}
 \Pr(\hbox{cycle})=(4/5)^3=64/125=0.512.
 \label{closure30:cycle-probability}
\end{equation}
Topological sorting detects this cycle; it does not make the graph acyclic.
If cyclic draws are discarded, the marginal probability of any one retained
arrow becomes
\begin{equation}
 \Pr(A\to B\mid\hbox{acyclic})
 =\frac{p(1-p^2)}{1-p^3}
 =\frac{36}{61}\simeq0.59016,\qquad p=4/5.
 \label{closure30:cycle-conditioning}
\end{equation}
Thus rejection changes the probabilities. Deleting an arrow from each cycle
also changes the distribution and requires a specified rule.

More fundamentally, no distribution supported on DAGs can retain all three
marginals at $0.8$. Every acyclic realization contains at most two of these
arrows. Taking expectations therefore gives
\begin{equation}
 \Pr(A\to B)+\Pr(B\to C)+\Pr(C\to A)\leq2,
 \label{closure30:dag-marginals}
\end{equation}
whereas the proposed sum is $2.4$. Enumeration of all eight arrow subsets
verifies both the conditional probabilities and this bound. This construction
shows what a general graph-sampling design must handle. It does not assert
that the thesis's seven-variable example actually generated a cycle.
Its stated topological procedure checks whether all vertices are ordered;
the inspected description does not establish a general repair that preserves
the original arrow probabilities.
\citep[pp.~49--50]{closure30molero}

\begin{figure}[tbp]
 \centering
 \includegraphics[width=0.96\textwidth]{twelve-hour-closure/2026-10-29/graph-figures/uncertainty.pdf}
 \caption{Three different uncertainty questions. Left: a success cap leaves
 a range of feasible probabilities. Centre: discarding cycles changes the
 marginal probability of each arrow in the constructed three-arrow example.
 Right: retaining two observationally compatible graphs adds zero to an
 interval around 0.6. These are project constructions, not measured bank or
 language-model performance.}
 \label{closure30:uncertainty-figure}
\end{figure}

\section{Agreement between fitted models is not coverage of an effect}

Take two centred Gaussian structural models, with independent disturbances
within each model:
\begin{align}
 G_1:&\quad A=U_A,\quad Y=0.6A+U_Y,
 &\operatorname{Var}(U_A)=1,\quad\operatorname{Var}(U_Y)=0.64;\nonumber\\
 G_2:&\quad Y=V_Y,\quad A=0.6Y+V_A,
 &\operatorname{Var}(V_Y)=1,\quad\operatorname{Var}(V_A)=0.64.
 \label{closure30:two-gaussian-worlds}
\end{align}
Both generate the same observational covariance,
\begin{equation}
 \Sigma=\begin{pmatrix}1&0.6\\0.6&1\end{pmatrix},
 \label{closure30:shared-covariance}
\end{equation}
but the effect of increasing $A$ by one unit is $0.6$ under $G_1$ and zero
under $G_2$. This is a deliberately static example with unresolved ordering,
such as two variables recorded together in a business survey. It does not
permit reversing an action and an outcome whose temporal order is known.
A verified information clock could exclude $G_2$; possible hidden common
causes would then raise a different identification question.

If both an estimated model and its reference impose $G_1$, they can agree
perfectly about the effect while $G_2$ generates the observations. In
particular, two identical nondegenerate intervals have overlap score one
whether or not either contains the true effect. For intervals $I$ and $J$,
the source's one-dimensional overlap score is
\begin{equation}
 \frac{\max\{0,\min(u_I,u_J)-\max(\ell_I,\ell_J)\}}
 {\max(u_I,u_J)-\min(\ell_I,\ell_J)}.
 \label{closure30:interval-agreement}
\end{equation}
The identical-interval value is an identity, not a new empirical result.
An overlap score remains useful for comparing models against a declared
reference; it must be reported as that comparison.
\citep[p.~61, equations (23)--(24)]{closure30molero}

A small repeated-sample study makes the distinction visible without neural
training. Generate 2,000 independent datasets of 500 observations from
\eqref{closure30:shared-covariance}. In each, regress $Y$ on an intercept and
$A$, and form the ordinary two-sided 95\% Student interval for its slope.
The Gaussian linear model makes this a valid slope interval. It is also a
valid intervention interval under $G_1$, but the same numerical calculation
does not become an intervention interval under $G_2$.

\begin{table}[tbp]
\centering
\caption{One observation law, two causal interpretations. Project simulation;
 2,000 repetitions, 500 observations per repetition.}
\begin{tabular}{lrr}
\toprule
Returned set & Coverage under $G_1$ & Coverage under $G_2$\\
\midrule
Regression interval $I_1$ & 95.7\% & 0.0\%\\
Graph-preserving set $I_1\cup\{0\}$ & 95.7\% & 100.0\%\\
\bottomrule
\end{tabular}
\label{closure30:graph-coverage-table}
\end{table}

The regression intervals have mean width 0.1411. The Monte Carlo standard
error of the observed 95.7\% coverage is about 0.45 percentage points; this
sample does not establish an improvement over the nominal 95\%. All 2,000
replication results are retained, and first, middle and last calculations
were independently checked by QR decomposition. The graph-preserving set
contains zero by construction and preserves the first world's interval.
It is a confidence set for the declared two-model class, not an identified
set over all latent or nonlinear mechanisms.

For an economic illustration, suppose one standardized outcome unit is worth
USD 100,000 and the one-unit change in $A$ costs USD 20,000. These are
transparent scenario inputs, not observed prices. The two stipulated effects
imply net values of USD 40,000 and minus USD 20,000. The ambiguity changes
the investment decision. Reporting only the narrow positive-effect interval
would conceal the reason for obtaining assignment or timing information
before committing resources.

\section{A constructive route from graph evidence to decision uncertainty}

The proposal can retain a set of admissible diagrams rather than requiring a
probability for every arrow. Let $\widehat{\mathcal G}$ be the retained class
and let $G_\star$ denote the true member of a declared scientific model class.
Suppose the procedure establishes
\begin{equation}
 \Pr\{G_\star\notin\widehat{\mathcal G}\}\leq\alpha_G.
 \label{closure30:graph-retention-contract}
\end{equation}
For each graph, obtain an identified effect set or an inferential set
$C_g$ whose relevant coverage event is valid jointly with the selection
procedure. In particular, suppose
$\Pr\{\theta_\star\notin C_{G_\star}\}\leq\alpha_E$.
Then return
\begin{equation}
 C=\bigcup_{g\in\widehat{\mathcal G}} C_g,
 \qquad
 \Pr\{\theta_\star\notin C\}\leq\alpha_G+\alpha_E.
 \label{closure30:union-coverage}
\end{equation}
To prove this, if the true graph is retained and its own set covers the
effect, their union covers it. Failure is contained in the union of those
two failure events. No independence between them is required.

This elementary guarantee identifies the work that a prototype must do.
An interval justified only for a graph fixed before looking at the data
does not automatically satisfy the effect condition after graph selection.
One route uses a simultaneous family of intervals over a fixed candidate
class; another uses an independent sample for effect estimation, subject to
the relevant assumptions. Sample splitting does not repair a false causal
model. If a retained graph does not identify the effect, its $C_g$ must
represent the justified bounds or remain uninformative. Dropping that graph
because it widens the answer would spend unaccounted graph error.

In the two-world calculation, the class is fixed and contains both graphs,
so $\alpha_G=0$. The reverse graph gives the exact set $\{0\}$; the forward
graph gives the ordinary slope interval. This explains the second row of
Table~\ref{closure30:graph-coverage-table}. In a real application the hard
part is establishing the retention and estimation conditions. A language
score alone supplies neither. The method is therefore a concrete proposal
route with explicit research obligations, not a claim that general calibrated
language-assisted causal inference has been solved.

\section{What to adopt, what to test and what to decline}

The public repository linked by the thesis is a PyWhy-LLM adaptation. Its
2 February 2026 revision contains the language suggestion and retrieval
components, pairwise notebooks, declared dependencies and an MIT notice
carrying the PyWhy copyright. The inspected complete tree and selected
notebooks do not contain the thesis's HallBayes integration, Monte Carlo
graph construction or causal-flow training pipeline. They support reuse
assessment of the front end, not independent reproduction of all reported
results. Its retrieval helper also creates its own language-model client and
falls back to an answer without retrieved evidence when no match is found.
A bank adapter should pass the selected model explicitly and retain the
absence of evidence as an observable condition. The proposal should
integrate the front-end ideas through a small,
versioned interface and prefer maintained upstream components where their
scope matches; porting the incomplete repository wholesale offers no
demonstrated advantage. Preserve applicable notices and review dependencies
and data permissions separately.
\citep{closure30molerocode}

The HallBayes repository at the latest inspected pre-cutoff revision presents
Berry, an evidence-support service with command-line and MCP interfaces.
Its README and source distinguish cited-evidence support from world truth.
That distinction fits a useful role in checking an assistant's attributed
claims. The June 2026 revision is not the December 2025 implementation cited
by the thesis. A date-bounded query of the current default-branch history
returned no commit for that earlier date, while a retained version-one tag
points to February 2026. This does not prove the older software never
existed; it prevents silently treating the current service as the thesis's
exact dependency. No downloaded source was executed in this comparison.
\citep{closure30berry}

The decision is consequently specific. Adopt the separation between
evidence retrieval, graph construction and intervention calculation.
Evaluate quotation and evidence-support checks for their own stated
predicates. Decline the direct conversion of an information ceiling into
calibrated edge probabilities, and decline unqualified confidence-interval
language for the spread of fitted-model effects. Before a prototype makes
stronger claims, it needs an authentic complete implementation, a declared
joint graph law or retention procedure, the relevant calibration evidence,
and public or synthetic experiments with known interventions and matched
information across comparators. Bank validation can follow when those
experiments justify its cost.

\section{A survey source that cannot support the proposed comparison}

A separate institutional thesis, \emph{Large Language Models for Causal
Inference}, was recovered from the University of Alcal\'a. Its identity is
clear, but its survey and productivity estimates cannot presently support
the proposal. It reports both ten and eighty included studies in different
places. Two arXiv identifiers assigned to causal-language studies in
Appendix A resolve to papers on Toeplitz determinants and a smart-building
security testbed. The named reproducibility repository currently returns
HTTP 404. These are specific, independently checkable discrepancies;
they do not establish intent or prove that a repository never existed.
\citep[Section 4.1, Tables 5--6 and 9, Section 5.4.3, Appendix A]
{closure30aitlhouss}

The source is retained for identity and exclusion transparency. Its reported
accuracy, time savings and independent trials are not used as evidence.
An authentic corrected study inventory and the underlying trial records
could change that disposition. This exclusion also illustrates a practical
requirement for a causal research assistant: locating a plausible citation
is not enough. The linked original must support the particular claim before
that claim enters a commercial comparison or an investment case.
